% Static vs dynamic equilibrium.
\begin{guidefig}{Static equilibrium: nothing moves. Dynamic equilibrium: things keep moving, but inputs and outputs balance, so the level stays the same, like body temperature or body weight.}
\begin{tikzpicture}
  % static: book on a table
  \node[font=\sffamily\bfseries,text=black!70] at (-3.8,2.1) {Static equilibrium};
  \fill[black!45] (-5.4,0) rectangle (-2.2,0.18); \fill[black!45] (-5.2,0) rectangle (-5,-1.1); \fill[black!45] (-2.6,0) rectangle (-2.4,-1.1);
  \fill[vocabc!70] (-4.4,0.18) rectangle (-3.2,0.55);
  \node[font=\scriptsize\sffamily,align=center] at (-3.8,1.2) {a book resting on a table:\\nothing moves};
  % dynamic: bathtub
  \node[font=\sffamily\bfseries,text=tipc] at (3.2,2.1) {Dynamic equilibrium};
  \fill[vocabc!25] (1.6,-1) rectangle (4.8,0.2);
  \draw[black!65,line width=1.4pt] (1.6,0.9) -- (1.6,-1) -- (4.8,-1) -- (4.8,0.9);
  \draw[black!40,dashed] (1.4,0.2) -- (5,0.2) node[right,font=\scriptsize\sffamily,text=black!70,align=left] {level stays\\the same};
  \draw[garr=vocabc,line width=1.6pt] (2.4,1.6) node[above,font=\scriptsize\sffamily,text=vocabc]{water in} -- (2.4,0.45);
  \draw[garr=vocabc,line width=1.6pt] (4.3,-1) -- (4.3,-1.7) node[below,font=\scriptsize\sffamily,text=vocabc]{water out};
  \node[font=\scriptsize\sffamily,align=center] at (3.2,-2.4) {in = out, so the level holds steady};
\end{tikzpicture}
\end{guidefig}
